\documentclass[11pt]{article}
\usepackage[T1]{fontenc}
\usepackage{textcomp}
\usepackage[margin=0.75in]{geometry}
\usepackage{amsmath,amssymb,amsthm}
\usepackage{array,booktabs}
\usepackage{hyperref}
\setlength{\parindent}{0pt}
\newtheorem{theorem}{Theorem}
\theoremstyle{definition}
\newtheorem{definition}{Definition}
\title{Flow Proof Artifact: prop-05-remainder-same-multiple}
\author{Generated by \texttt{flow doc proof}}
\date{Flow Proof Book}
\begin{document}
\maketitle
If one magnitude is the same multiple of another that a part subtracted is of a part subtracted, the remainder is the same multiple of the remainder.\\[0.5em]
\textbf{Source.} euclid — Elements, Book V, Proposition 5\\[0.5em]
\bigskip
\noindent{\large \textbf{Derived fact 1.} \textit{If one magnitude is the same multiple of another that a part subtracted is of a part subtracted, the remainder is the same multiple of the remainder} \hfill the Euclidean plane \cdot Euclid Book V \cdot Proposition 5: the remainder is the same multiple of the remainder}\par
\smallskip\noindent\textit{Coordinate: the Euclidean plane \cdot Euclid Book V \cdot Proposition 5: the remainder is the same multiple of the remainder}\par
\smallskip\noindent\textit{Source: euclid — Elements, Book V, Proposition 5}\par
\smallskip\noindent\textit{Built on: proposition 3: the difference of equimultiples is an equimultiple of the difference, for Euclid Book V on the Euclidean plane}\par
\medskip
\noindent\textbf{Goal.} If one magnitude is the same multiple of another that a part subtracted is of a part subtracted, the remainder is the same multiple of the remainder\par
\noindent$\displaystyle \text{remainder same multiple}$\par
\medskip
\noindent
\begin{tabular}{@{} >{\raggedright\arraybackslash}p{0.06\textwidth} >{\raggedright\arraybackslash}p{0.47\textwidth} >{\raggedleft\arraybackslash}p{0.41\textwidth} @{}}
\textbf{\#} & \textbf{Proof} & \textbf{Mathematics} \\
\midrule
\textbf{1.} & We prove that proposition 5: the remainder is the same multiple of the remainder for Euclid Book V the Euclidean plane. & \\[0.45em]
\textbf{2.} & Let whole = first\_magnitude. & \\[0.45em]
\textbf{3.} & Let part = subtracted\_part. & \\[0.45em]
\textbf{4.} & Let remainder = whole - part. & \\[0.45em]
\textbf{5.} & From step 2, step 3, and step 4, we can deduce that equimultiple remainder holds. & $\displaystyle \text{equimultiple remainder holds}$ \\[0.45em]
\textbf{6.} & We invoke the derived fact governing Euclid Book V the Euclidean plane: proposition 3: the difference of equimultiples is an equimultiple of the difference, for Euclid Book V on the Euclidean plane. & \\[0.45em]
\textbf{7.} & From step 6, this implies remainder same multiple. Hence proven. & $\displaystyle \text{remainder same multiple}$ \\[0.45em]
\bottomrule
\end{tabular}
\label{thm:Geometry-Euclid Book V-proposition_5_the_remainder_is_the_same_multiple_of_the_remainder}
\par\medskip\hrule
\end{document}
